\subsection{Exploratory Multi-Family Pilot: Search Survives, Selection Does Not}
\label{sec:pilot-results}

The single-event pilot tests a different bottleneck from the load-support bank: can the true physical explanation remain in a tractable four-family shortlist when onset and severity are fitted? The author-supplied campaign summaries report that a development-selected, frozen screening and racing policy retained the true candidate in all four blind mini-holdout cases (one fault, load change, generator change, and line outage). The generator case reached the six-survivor cutoff with no spare place. Four quiet windows provided only a pilot check of proposal scores; they cannot establish operational false-alarm control. The solver diagnostics also report full-block/Schur agreement on 30 synthetic systems and fault-resistance estimates of 0.02006--0.02133 from four initializations around a synthetic truth of 0.021. These checks establish numerical behavior of components on the reported cases, not a validated source-localization rate.

The same dossiers report a negative strict-final outcome: none of the four truth candidates won the final comparison. That 0/4 is currently a terminal-reported result without its matching final-survivor archive in the repository. Its mechanisms must remain distinct. Fault and line truth/competitor fits did not all reach strict convergence, so their ranking is numerically unresolved. In the load and generator cases, converged neighboring same-family candidates reportedly received better final scores. That pattern motivates an oracle sequence for event-operator parity, pre-event state, onset and duration, severity, noise, and counterfactual PMU coverage. It does not yet prove an intrinsic identifiability limit. The pilot thus supports continued work on the joint numerical route but supplies no positive multi-family end-to-end claim. Supplementary Appendix~H gives the case ledger, reported diagnostics, and the exact artifact checks needed to turn this audit into primary evidence.
